% OWNER: Nathan. Target: about 0.75 page including diagram/configuration.
\subsection{Systems and Controlled Inference}
\label{sec:systems}
\placeholder{Nathan}{Describe SeamlessM4T v2-large Direct ST and Whisper
large-v2 followed by NLLB-600M or NLLB-3.3B. Cite the v2 context, Whisper and
NLLB sources~\cite{seamless2023,radford2023whisper,nllb2022}. Give exact checkpoint
IDs/resolved revisions and target/source language codes. Explain that model
families, training resources and objectives differ.}

\placeholder{Nathan}{Specify mono 16 kHz audio, actual fixed greedy settings,
seed, FP16/T4 environment and relevant library versions from run records.
Explain staged ASR then MT and no reference transcripts supplied to MT. Describe
pilots as engineering checks, checkpointing, successful-run/limit checks, ID and
fingerprint alignment, and matching archived ASR strings from separate runs.
If runtime is reported, distinguish Direct pipeline time, Cascade stage-sum time,
and wall-clock time; loading/warm-up are excluded from recorded sample timings.}

\begin{figure}[t]
\centering
\figureplaceholder{Nathan: matched English speech, Direct ST, Whisper followed
by both NLLB variants, and common Chinese reference evaluation.}
\caption{System comparison. TODO (Nathan): replace with an original or properly
credited diagram; adapt the survey diagram to the completed three-system design.}
\label{fig:systems}
\end{figure}

% Optional compact table: include essential settings in prose if a table is wider
% or less readable. Keep full model revisions in the paper if needed for replay.
